\section{The literal native-multiplier interface}
\label{sec:native-interface}

\subsection{The frequency-free branch key}

The resolved TPC-93 key
\(\xi=(\theta,c,\kappa,r)\) contains the low frequency \(r\).
It must not be inserted into a census and then paired with every
frequency a second time.  We therefore use the frequency-free
TPC-94 branch
\[
 \beta=(\theta,c,\kappa).
\]
Let \(\mathfrak B_X^{\mathrm{nc}}\) be the finite set of soluble,
nonempty literal branches \(\beta\) for which all inherited masks
hold,
\[
 c\mid\Omega_\beta,\qquad q_X\nmid\Omega_\beta.
\]
Throughout this section every sum over \(\beta\) is restricted to
\(\mathfrak B_X^{\mathrm{nc}}\).
Here \(\theta\) retains the polarization, opposite row,
\(\ell,j,\sigma,v\), interval component, actual masks, projector
labels, and all continuous separation data.  It fixes the reflected
ultra family, not one value of the opened divisor: the latter is
\(U_\theta(t)\) along the affine fiber.

For every soluble, nonempty branch put
\[
 b_\beta=c\kappa,\qquad
 N_\beta=\#I_\beta^\ast,\qquad
 \Omega_\beta=\ell_\theta v_\theta\sigma_\theta B_\beta.
\]
Thus \(N_\beta\) is the cardinality of the resolved
\(z\)-fiber \(I_\beta^\ast\); it is not an orbit length in \(j\).
Let \(A_{\beta,X}\) be the exact TPC-94 outer coefficient and let
\begin{equation}
 \mathcal A_{\beta,X}
 =
 \sum_{z\in I_\beta^\ast}
 |\mathcal W_{\beta,X}(z)|
 \mu^2(\widetilde D_\beta(z))
 \mu^2(V_\beta(z)).
\label{eq:uniform-branch-mass}
\end{equation}
The exact branch sum satisfies
\[
 |S_{\beta,X}(r)|\le\mathcal A_{\beta,X}
 \quad\text{for every }r,
\]
so the frequency-independent proof-carrying weight is
\begin{equation}
 w_{\beta,X}
 =
 |A_{\beta,X}|\,\mathcal A_{\beta,X}.
\label{eq:base-branch-weight}
\end{equation}
This is the literal multiplier-mass return of TPC-94, not a new
uniformity assumption \citep{WangTPC94}.

On the nonconstant branch, literal provenance gives
\(c\mid\Omega_\beta\) and \(q_X\nmid\Omega_\beta\).  Define
\begin{equation}
 \omega_\beta
 =
 \varepsilon_\theta\frac{\Omega_\beta}{c}
 \pmod{q_X}
 \in\mathbb F_{q_X}^{\times}.
\label{eq:native-omega}
\end{equation}
The sign remains in the source key even though \(E_H=-E_H\).
Branches with \(N_\beta>q_X\) have no resonant nonzero frequency by
\cref{lem:exact-resonance-window}; they contribute zero to the
carrier below and are not discarded from the full signed formula.

For each exact width \(H\), define
\begin{equation}
 a_H(\omega)
 =
 \sum_{\substack{\beta:\ 1\le N_\beta\le q_X\\
  H_{q_X}(N_\beta)=H\\
  \omega_\beta=\omega}}
 w_{\beta,X}.
\label{eq:native-census}
\end{equation}
Equal multipliers are grouped only after every branch has received
its own weight; the provenance bag remains recoverable.

Let
\[
 b_X(r)=|\mathfrak m_{K,X}(r)|
 \one_{\{0<d_{q_X}(r,0)\le R_X\}},
\]
extended by zero to \(G\).  TPC-93 gives
\begin{equation}
 \|b_X\|_1\ll_K1,\qquad
 \|b_X\|_2\ll_K1.
\label{eq:filter-norms}
\end{equation}

\begin{definition}[Positive literal resonance carrier]
\label{def:positive-carrier}
Set
\begin{equation}
 \mathcal C_{\mathrm{res},X}^{\mathrm{abs}}
 =
 \sum_{\substack{\beta:\ 1\le N_\beta\le q_X}}
 w_{\beta,X}
 \sum_{r\in G}
 b_X(r)\,
 \one_{E_{H_{q_X}(N_\beta)}}(r\omega_\beta).
\label{eq:positive-carrier}
\end{equation}
\end{definition}

\begin{theorem}[Literal positive majorant and exact census identity]
\label{thm:literal-carrier}
The absolute value of the signed nonconstant resonant contribution
is at most \(\mathcal C_{\mathrm{res},X}^{\mathrm{abs}}\), and
\begin{equation}
 \boxed{
 \mathcal C_{\mathrm{res},X}^{\mathrm{abs}}
 =
 \sum_{H=1}^{(q_X-1)/2}
 \mathcal I_H(a_H,b_X).}
\label{eq:carrier-sum}
\end{equation}
For each \(H\),
\begin{align}
 \mathcal I_H(a_H,b_X)
 ={}&
 \frac{2H}{q_X-1}
 \|a_H\|_1\|b_X\|_1
 +\mathcal R_H(a_H,b_X),
\label{eq:literal-principal-remainder}\\
 |\mathcal R_H(a_H,b_X)|
 \ll_K{}&
 \sqrt{q_X}\log q_X\,\|a_H^\circ\|_2.
\label{eq:literal-remainder-bound}
\end{align}
\end{theorem}

\begin{proof}
In the exact TPC-94 sector formula, the outer coefficient and branch
mass obey
\[
 |A_{\beta,X}\mathfrak m_{K,X}(r)
   \zeta_{\beta,r}S_{\beta,X}(r)|
 \le w_{\beta,X}b_X(r).
\]
The unit phase has disappeared only after taking absolute values.
\Cref{lem:exact-resonance-window} supplies the indicator in
\eqref{eq:positive-carrier}; for \(N_\beta>q_X\) it is identically
zero.  This proves the majorization without pairing one branch with
any unintended frequency.  Grouping the finite sum by the exact
pair \((H,\omega_\beta)\) gives the identity
\eqref{eq:carrier-sum}.  Now apply
\cref{thm:principal-nonprincipal,cor:spectral-bound} and
\eqref{eq:filter-norms}.
\end{proof}

\begin{corollary}[A concrete sufficient gate]
\label{cor:sufficient-gate}
The complete resonance carrier is
\(O(X^{o(1)}Q_X^2)\) if
\begin{align}
 \sum_H\frac{H}{q_X}\|a_H\|_1
 &\ll X^{o(1)}Q_X^2,
\label{eq:principal-gate}\\
 \sum_H\|a_H^\circ\|_2
 &\ll
 \frac{X^{o(1)}Q_X^2}
 {\sqrt{q_X}\log q_X}.
\label{eq:dispersion-gate}
\end{align}
Since \(q_X\asymp Q_X\), the second target is
\[
 \sum_H\|a_H^\circ\|_2
 \ll
 X^{o(1)}Q_X^{3/2}.
\]
\end{corollary}

\begin{proof}
Sum \eqref{eq:literal-principal-remainder} over \(H\), use
\eqref{eq:filter-norms}, and insert
\eqref{eq:principal-gate}--\eqref{eq:dispersion-gate}.
\end{proof}

\begin{remark}[Sufficient, not necessary]
The absolute carrier discards signed cancellation inside each affine
sum and across native keys.  Failure of
\eqref{eq:principal-gate} or \eqref{eq:dispersion-gate} is not a
lower bound for the carrier and does not disprove the original
signed route.  A proof-carrying polynomial obstruction in one of
the actual termwise bounds only stops this sufficient
principal-plus-absolute-remainder implementation; a sharper
filter-sensitive estimate could retain cancellation between those
pieces.
\end{remark}
